\chapter{Spontaneous Supersymmetry Breaking and the First \texorpdfstring{\(SU(5)\)}{SU(5)} Multiplet}

\lectureoverview{Symmetry breaking first appears here as a visible fact about spectra: a rotationally invariant atom has degenerate angular-momentum levels, while a magnetic field splits them. A ferromagnet then teaches the subtler lesson that the equations may retain a symmetry which the vacuum does not. Gapless magnons, the condition \(J|0\rangle\neq0\), and a nonzero order parameter are three descriptions of that same phenomenon. We carry this pattern into a chiral supersymmetric field theory, derive the scalar potential \(V=|W'|^2\) by eliminating the auxiliary field \(F\), and distinguish supersymmetric from \(F\)-breaking vacua. The resulting vacuum energy, boson--fermion splitting, goldstino, hidden sector, and super-Higgs mechanism complete the supersymmetry discussion. The final part changes subject and opens grand unification: \(SU(3)\times SU(2)\times U(1)\) is embedded blockwise in \(SU(5)\), whose first matter multiplet combines a lepton doublet with three anti-down quarks.}

\section{A Spectrum Remembers Its Symmetry}

The purpose of beginning with an atom is not merely to review elementary quantum mechanics. Energy levels and particle masses are the same kind of spectral data. If a symmetry forces atomic levels to coincide, then its breaking can split those levels; in field theory the corresponding statement is that members of a particle multiplet can acquire different masses.

Let an atomic state have total angular momentum \(j\), including both spin and orbital contributions. Relative to any chosen axis,
\begin{equation}
J_z|j,m\rangle=m\hbar|j,m\rangle,
\qquad
m=-j,-j+1,\ldots,j,
\end{equation}
so there are \(2j+1\) states. The familiar textbook arrows should not be read as saying that the angular momentum vector literally has only \(2j+1\) possible spatial directions. What is quantized is its component along the chosen axis.

In empty, rotationally invariant space there is no physical reason to prefer that axis. A rotation mixes the states within the spin-\(j\) multiplet, and invariance of the Hamiltonian implies
\begin{equation}
[H,\mathbf J]=0,
\qquad
E_{j,-j}=E_{j,-j+1}=\cdots=E_{j,j}.
\label{eq:l9-degeneracy}
\end{equation}

\begin{classroomqa}
\audiencequestion{What symmetry says that rotating the whole atom cannot change its energy? \lecturetimestamp{00:05:53}}
\lecturerresponse{Rotational symmetry. Since a rotation carries one orientation of the angular momentum into another without changing the physical system, all \(2j+1\) states have the same energy.}
\end{classroomqa}

Now apply a magnetic field \(\mathbf B=B\hat{\mathbf z}\). The interaction
\begin{equation}
H_B=-\boldsymbol{\mu}\cdot\mathbf B
\end{equation}
distinguishes the component parallel to the field. Schematically,
\begin{equation}
\Delta E_m\propto mB.
\label{eq:l9-zeeman}
\end{equation}
Levels with different \(m\) separate linearly at weak field, and transitions among them split spectral lines.

\begin{classroomqa}
\audiencequestion{What is the splitting of atomic levels in a magnetic field called? \lecturetimestamp{00:07:13}}
\lecturerresponse{The Zeeman effect. It is the prototype for the entire discussion: levels made degenerate by a symmetry separate when a background breaks that symmetry.}
\end{classroomqa}

\begin{center}
\begin{tikzpicture}[x=1.05cm,y=0.75cm,>=Latex]
  \draw[thick] (-3,0) -- (-1,0);
  \node[left] at (-3,0) {\(B=0\)};
  \foreach \y/\m in {-1.2/-1,-0.4/0,0.4/1,1.2/2} {
    \draw[thick,teal!70!black] (1,\y) -- (3,\y);
    \node[right] at (3,\y) {\(m=\m\)};
  }
  \draw[->,thick] (-0.5,0) -- (0.6,0);
  \node[above] at (0.05,0) {\(\mathbf B\neq0\)};
\end{tikzpicture}
\end{center}

\section{The Ferromagnet: Symmetric Laws, Asymmetric Vacuum}

A simple microscopic picture of a ferromagnet assigns a small magnetic moment to each site. If neighboring moments lower their energy by becoming parallel, the low-temperature phase is ferromagnetic. If antiparallel alignment is preferred, the result is antiferromagnetic. Strong thermal or quantum fluctuations can instead wash out any net magnetization.

Suppose a ferromagnet fills the world and settles into a uniformly aligned state. Its internal magnetic field splits the levels of an atom just as an externally imposed field would. Yet there is a decisive difference. The underlying interactions do not prefer north, east, or any other direction. Every uniformly rotated configuration has the same energy:
\begin{equation}
\mathbf M(\mathbf x)=M_0\hat{\mathbf n},
\qquad
\hat{\mathbf n}\in S^2.
\label{eq:l9-ferromagnetic-vacua}
\end{equation}
The equations are rotationally invariant, but one actual ground state chooses one \(\hat{\mathbf n}\). This is spontaneous rather than explicit symmetry breaking.

An inhabitant made from excitations of this medium cannot grab the entire universe and rotate it. Locally, therefore, the chosen direction looks as fundamental as an external field. An atom still sees the same level pattern, with its angular momentum resolved along \(\hat{\mathbf n}\) rather than along an externally named \(z\)-axis. The question is how an internal observer could discover that the apparent asymmetry belongs to the state, not to the laws.

\subsection{The Goldstone Diagnostic}

A uniform rotation of every moment moves from one vacuum in Eq.~\eqref{eq:l9-ferromagnetic-vacua} to another and costs no energy. Let the orientation instead change very gradually across space. Neighboring moments then pull on one another and the distortion propagates as a spin wave, whose quanta are magnons. As the wavelength tends to infinity, the configuration approaches a uniform rotation, so its excitation energy tends to zero:
\begin{equation}
k=\frac{2\pi}{\lambda}\longrightarrow0
\quad\Longrightarrow\quad
E(k)\longrightarrow0.
\label{eq:l9-gapless-magnon}
\end{equation}
Two inhabitants can wiggle nearby moments, send such waves across the medium, and measure this arbitrarily low energy. That experiment reveals spontaneous breaking even though they cannot rotate the world.

\editorialnote{A ferromagnetic magnon is gapless but generally has a nonrelativistic quadratic dispersion \(E\propto k^2\), rather than the relativistic relation \(E\propto|k|\). The lecture uses the word massless in the physically relevant sense of vanishing gap. Both are Goldstone modes.}

The same idea applies to continuous internal symmetries, including the gauge and flavor groups of particle physics. When the broken symmetry is gauged, however, the Goldstone mode does not remain as an independent massless particle.

\begin{classroomqa}
\audiencequestion{What happens when spontaneous symmetry breaking is combined with gauge invariance? \lecturetimestamp{00:15:56}}
\lecturerresponse{The Goldstone boson supplies an additional polarization to the gauge field. The gauge boson becomes massive; this is the Higgs mechanism, to which we will return in the gravitational analogy.}
\end{classroomqa}

\section{Three Equivalent Signals of Breaking}

The ferromagnet can now be described in three mutually reinforcing languages.

\paragraph{Degenerate vacua.}
There is one minimum-energy state for every orientation \(\hat{\mathbf n}\). Physics takes place in one chosen member of this family even though the symmetry relates all of them.

\paragraph{A generator that does not annihilate the vacuum.}
Let \(J_a\) be a generator of rotations. A finite transformation is
\begin{equation}
U(\epsilon)=e^{-i\epsilon J_a},
\end{equation}
and \(J_a|0\rangle\) is proportional to the infinitesimal change of the vacuum, not the entire rotated state.

\begin{classroomqa}
\audiencequestion{What are the generators of rotational symmetry? \lecturetimestamp{00:17:29}}
\lecturerresponse{The components of total angular momentum. We call them \(J_a\) because they rotate everything in the state, not merely one atom.}
\end{classroomqa}

An unbroken generator satisfies
\begin{equation}
J_a|0\rangle=0.
\label{eq:l9-unbroken-generator}
\end{equation}
A broken generator instead gives
\begin{equation}
J_a|0\rangle\neq0.
\label{eq:l9-broken-generator}
\end{equation}

\begin{classroomqa}
\audiencequestion{If \(J_a|0\rangle\) is nonzero, does it simply give the same vacuum? \lecturetimestamp{00:19:43}}
\lecturerresponse{No. It gives the infinitesimal difference between the chosen vacuum and a slightly rotated one. Physically that difference is an infinitely long-wavelength Goldstone excitation.}
\end{classroomqa}

\begin{classroomqa}
\audiencequestion{Does this infinitesimal rotation change the energy? \lecturetimestamp{00:20:14}}
\lecturerresponse{No. The rotated vacuum is degenerate with the original one, so the zero-momentum Goldstone state has zero energy.}
\end{classroomqa}

Symbolically,
\begin{equation}
J_a|0\rangle
\sim
|\pi_a,\mathbf k=0\rangle,
\label{eq:l9-generator-goldstone}
\end{equation}
where \(\pi_a\) denotes the Goldstone mode associated with the broken generator.

\paragraph{A nonzero order parameter.}
The local magnetization \(\mathbf M(\mathbf x)\) is a field. In a warm or quantum-disordered phase,
\begin{equation}
\langle\mathbf M\rangle=0,
\end{equation}
and the vacuum does not select a direction. In the ordered phase,
\begin{equation}
\langle M_i\rangle\neq0
\quad\text{for at least one component }i,
\label{eq:l9-magnetization-order}
\end{equation}
so a rotation mixes that nonzero component into the others. Degenerate vacua, a broken generator, and a nonzero order parameter are the three patterns to be carried into supersymmetry.

\section{The Chiral Multiplet Recalled}

The simplest four-dimensional example contains one complex scalar, one Weyl fermion, and one auxiliary complex scalar. In chiral coordinates,
\begin{equation}
y^\mu=x^\mu+i\bar\theta\bar\sigma^\mu\theta,
\end{equation}
a chiral superfield may be expanded schematically as
\begin{equation}
\Phi(y,\theta)=
\phi(y)+\sqrt2\,\theta\psi(y)+\theta^2F(y).
\label{eq:l9-chiral-expansion}
\end{equation}
The lecture suppresses the conventional factor \(\sqrt2\); we include it from this point so that the component coefficients are standard. A chiral superfield is intrinsically complex: imposing reality is not preserved by its supersymmetry transformation.

The supersymmetric kinetic action is
\begin{equation}
S_{\mathrm{kin}}
=
\int d^4x\,d^4\theta\,
\Phi^\dagger\Phi.
\label{eq:l9-kahler-action}
\end{equation}
Grassmann integration selects the coefficient of
\(\theta^2\bar\theta^2\). With mostly-minus metric, the component result is
\begin{equation}
\mathcal L_{\mathrm{kin}}
=
\partial_\mu\phi^\ast\,\partial^\mu\phi
+i\bar\psi\bar\sigma^\mu\partial_\mu\psi
+F^\ast F.
\label{eq:l9-chiral-kinetic}
\end{equation}
The first two terms propagate the scalar and fermion. The last has no derivative and therefore propagates no particle.

\begin{figure}[H]
\centering
\includegraphics[width=0.76\textwidth]{lecture_09_figure_02.png}
\lectureframe{00:27:36.000}
\caption{The chiral coordinate and component expansion above the scalar, fermion, and auxiliary terms.}
\end{figure}

\section{What the \texorpdfstring{\(d^2\theta\)}{d2-theta} Integral Selects}

Interactions and masses arise from a holomorphic function \(W(\Phi)\), the superpotential:
\begin{equation}
S_W
=
\int d^4x
\left[
\int d^2\theta\,W(\Phi)
+\mathrm{h.c.}
\right].
\label{eq:l9-superpotential-action}
\end{equation}
A chiral expression has no independent \(\bar\theta\) dependence, so integrating it over all four odd coordinates would give zero. The \(d^2\theta\) integral instead selects its \(\theta^2\) component.

To see the combinatorics before taking a general function, let \(W(\Phi)=\Phi^n\). In
\begin{equation}
(\phi+\sqrt2\,\theta\psi+\theta^2F)^n,
\end{equation}
there are two ways to obtain exactly two \(\theta\)'s. One may choose a single \(F\theta^2\), or choose two fermionic factors.

\begin{classroomqa}
\audiencequestion{How many terms contain one chosen \(F\theta^2\) and otherwise only scalars: \(n-1\) or \(n\)? \lecturetimestamp{00:31:33}}
\lecturerresponse{There are \(n\). Any one of the \(n\) factors may supply \(F\theta^2\); the remaining \(n-1\) factors supply \(\phi\). The contribution is \(nF\phi^{n-1}\).}
\end{classroomqa}

\begin{classroomqa}
\audiencequestion{How many ways can two factors supply the two fermions? \lecturetimestamp{00:32:26}}
\lecturerresponse{There are \(\binom n2=n(n-1)/2\) unordered choices. Grassmann contraction fixes the corresponding \(\psi\psi\) term.}
\end{classroomqa}

Thus, in the lecture's coefficient-suppressing notation,
\begin{equation}
\int d^2\theta\,\Phi^n+\mathrm{h.c.}
\;\leadsto\;
nF\phi^{n-1}
-\frac{n(n-1)}{2}\phi^{n-2}\psi\psi
+\mathrm{h.c.}
\label{eq:l9-power-extraction}
\end{equation}
The minus sign shown here follows the standard normalized expansion and Grassmann convention; changing conventions can move superficial signs.

\begin{classroomqa}
\audiencequestion{What happened to the term \(\phi^n\), which contains no \(\theta\)'s at all? \lecturetimestamp{00:33:10}}
\lecturerresponse{It is present in the product, but \(d^2\theta\) discards it. The arrow means that integration selects the \(\theta^2\) coefficient; it is not equality of the full expressions.}
\end{classroomqa}

\begin{figure}[H]
\centering
\includegraphics[width=0.82\textwidth]{lecture_09_figure_03.png}
\lectureframe{00:33:38.000}
\caption{The repeated superfield factors and the implication arrow that extracts, rather than equals, their two-\(\theta\) component.}
\end{figure}

Summing powers with arbitrary coefficients converts the combinatorial factors into derivatives:
\begin{equation}
\int d^2\theta\,W(\Phi)+\mathrm{h.c.}
\;\leadsto\;
F W'(\phi)
-\frac12 W''(\phi)\psi\psi
+\mathrm{h.c.}
\label{eq:l9-W-components}
\end{equation}

\begin{classroomqa}
\audiencequestion{Why does the fermion term contain \(\phi^{n-2}\), rather than \(\phi^{n-1}\)? \lecturetimestamp{00:36:32}}
\lecturerresponse{Two factors have already supplied the fermions, leaving \(n-2\) scalar factors. The coefficient is one half of the second derivative \(W''(\phi)\).}
\end{classroomqa}

The fermion term is important: it gives a mass when \(W''\) is constant and a Yukawa interaction when \(W''\) depends on \(\phi\). For the vacuum analysis, however, the auxiliary term \(FW'\) carries the immediate information.

\section{Eliminating the Auxiliary Field}

Collect the terms that contain \(F\):
\begin{equation}
\mathcal L_F
=
F^\ast F
+F W'(\phi)
+F^\ast \overline{W'(\phi)}.
\label{eq:l9-LF}
\end{equation}
Since there is no \(\partial_\mu F\), the full Euler--Lagrange equation
\begin{equation}
\partial_\mu
\frac{\partial\mathcal L}{\partial(\partial_\mu F)}
-
\frac{\partial\mathcal L}{\partial F}
=0
\end{equation}
reduces immediately to an algebraic equation. Treating \(F\) and \(F^\ast\) as independent variables gives
\begin{align}
\frac{\partial\mathcal L_F}{\partial F}
&=F^\ast+W'(\phi)=0,
&
\frac{\partial\mathcal L_F}{\partial F^\ast}
&=F+\overline{W'(\phi)}=0,
\label{eq:l9-F-eom}
\end{align}
and hence
\begin{equation}
F^\ast=-W'(\phi),
\qquad
F=-\overline{W'(\phi)}.
\label{eq:l9-F-solution}
\end{equation}

\begin{figure}[H]
\centering
\includegraphics[width=0.82\textwidth]{lecture_09_figure_04.png}
\lectureframe{00:39:40.000}
\caption{The complete auxiliary-field expression and its algebraic equation of motion, visible together after the board is cleared of occlusion.}
\end{figure}

\begin{classroomqa}
\audiencequestion{Should that derivative be taken with respect to \(\phi\), rather than the conjugate field written on the board? \lecturetimestamp{00:41:21}}
\lecturerresponse{Yes. The correction restores the conjugate pair in Eq.~\eqref{eq:l9-F-eom}: \(F^\ast=-W'(\phi)\) and \(F=-\overline{W'(\phi)}\).}
\end{classroomqa}

Substitution into Eq.~\eqref{eq:l9-LF} gives
\begin{equation}
\mathcal L_F\big|_{\mathrm{on\ shell}}
=-|W'(\phi)|^2.
\end{equation}
Since a scalar Lagrangian has the form kinetic energy minus potential energy, the ordinary scalar potential is
\begin{equation}
\boxed{V(\phi,\phi^\ast)=|W'(\phi)|^2\ge0.}
\label{eq:l9-F-potential}
\end{equation}
This \(V\) is not the superpotential. We reserve \(W\) for the holomorphic superspace function and \(V\) for ordinary potential energy.

\subsection{Two Immediate Examples}

With exact conventional normalizations,
\begin{equation}
\begin{gathered}
W(\Phi)=\frac12m\Phi^2,\\
W'(\phi)=m\phi,
\qquad V=|m|^2|\phi|^2.
\end{gathered}
\label{eq:l9-quadratic-W}
\end{equation}
The same term gives the fermion mass
\(-\tfrac12m\psi\psi+\mathrm{h.c.}\), so the scalar and fermion masses agree before supersymmetry is broken.

For
\begin{equation}
W(\Phi)=\frac{g}{3}\Phi^3,
\end{equation}
one obtains
\begin{equation}
\begin{gathered}
W'(\phi)=g\phi^2,
\qquad V=|g|^2|\phi|^4,\\
\mathcal L_{\mathrm{Yukawa}}
=-g\phi\psi\psi+\mathrm{h.c.}
\end{gathered}
\label{eq:l9-cubic-W}
\end{equation}
The quartic scalar vertex describes two-scalar scattering, while the Yukawa term is its supersymmetric companion.

\section{Two Kinds of Vacuum}

The vacuum is a constant field configuration at a minimum of \(V\); its derivative terms vanish. Because Eq.~\eqref{eq:l9-F-potential} is nonnegative, two cases are possible.

\subsection{A Zero-Energy Minimum}

If some \(\phi_0\) satisfies
\begin{equation}
W'(\phi_0)=0,
\end{equation}
then
\begin{equation}
V(\phi_0)=0,
\qquad
\langle F\rangle=0.
\label{eq:l9-unbroken-F}
\end{equation}
The vacuum is supersymmetric. This is the analogue of vanishing magnetization: no auxiliary component has selected a superspace direction.

\subsection{A Lifted Minimum}

If the equations
\begin{equation}
W'(\phi)=0
\label{eq:l9-F-flatness}
\end{equation}
have no simultaneous solution at the minimum, then
\begin{equation}
V_{\min}>0,
\qquad
\langle F\rangle\neq0,
\label{eq:l9-broken-F}
\end{equation}
and supersymmetry is spontaneously broken. The action remains supersymmetric; its lowest-energy state does not.

\begin{figure}[H]
\centering
\includegraphics[width=0.88\textwidth]{lecture_09_figure_05.png}
\lectureframe{00:52:09.960}
\caption{The board keeps the positive potential, the \(F\)-equation, and the vacuum condition in view while distinguishing symmetric equations from a non-symmetric solution.}
\end{figure}

The order-parameter statement can be made precise through the fermion transformation:
\begin{equation}
\delta\psi_\alpha
=
i\sqrt2
(\sigma^\mu\bar\epsilon)_\alpha\partial_\mu\phi
+\sqrt2\,\epsilon_\alpha F.
\label{eq:l9-fermion-transformation}
\end{equation}
For a homogeneous vacuum the derivative term vanishes. If
\(\langle F\rangle\neq0\), supersymmetry shifts the fermion by a nonzero constant, so the vacuum cannot be invariant.\editorialnote{Equation~\eqref{eq:l9-fermion-transformation} is the standard component form that makes explicit the lecture's statement that superfield components rotate into one another.}

The same conclusion follows from the superalgebra. In the rest frame,
\begin{equation}
\{Q_\alpha,Q_{\dot\alpha}^\dagger\}
=2\sigma^0_{\alpha\dot\alpha}H.
\end{equation}
Taking a vacuum expectation value expresses its energy as a sum of norms of supercharges acting on the vacuum:
\begin{equation}
E_0\ge0,
\qquad
E_0=0
\ \Longleftrightarrow\
Q_\alpha|0\rangle
=Q_{\dot\alpha}^\dagger|0\rangle=0.
\label{eq:l9-algebra-vacuum-energy}
\end{equation}
Thus positive vacuum energy and spontaneous supersymmetry breaking are two sides of the same rigid-supersymmetry statement.

\begin{center}
\begin{tikzpicture}[x=1.2cm,y=1.0cm,>=Latex]
  \draw[->] (-3,0) -- (-0.2,0) node[right] {\(\phi\)};
  \draw[->] (-2.8,-0.15) -- (-2.8,2.1) node[above] {\(V\)};
  \draw[thick,teal!75!black,domain=-2.5:-0.5,samples=60]
    plot (\x,{0.65*(\x+1.5)^2});
  \fill (-1.5,0) circle (2pt);
  \node[below] at (-1.5,0) {\(F=0\)};
  \draw[->] (0.4,0) -- (3.2,0) node[right] {\(\phi\)};
  \draw[->] (0.6,-0.15) -- (0.6,2.1) node[above] {\(V\)};
  \draw[thick,orange!85!black,domain=0.9:2.9,samples=60]
    plot (\x,{0.55*(\x-1.9)^2+0.65});
  \fill (1.9,0.65) circle (2pt);
  \draw[dashed] (0.6,0.65) -- (1.9,0.65);
  \node[below] at (1.9,0.65) {\(F\neq0\)};
\end{tikzpicture}
\end{center}

\editorialnote{A minimal illustration is \(W=f\Phi\), for which \(W'=f\) never vanishes and \(V=|f|^2\). That one-field example has a flat scalar direction; realistic stable \(F\)-term breaking is normally implemented with several chiral fields, as in O'Raifeartaigh-type hidden sectors.}

\section{What a Broken Supersymmetry Produces}

\subsection{Mass Splitting}

The ferromagnetic analogy now becomes spectral again. A weak internal magnetic field separates atomic levels without destroying the identity of the atom. Similarly, a weakly communicated supersymmetry-breaking order parameter separates bosons and fermions that would otherwise lie in one supermultiplet.

\begin{classroomqa}
\audiencequestion{Which particles are rotated into one another by supersymmetry and therefore lose their mass degeneracy when the vacuum breaks it? \lecturetimestamp{00:51:40}}
\lecturerresponse{Bosons and fermions. Their masses need not remain equal once the state no longer respects the transformation that relates them.}
\end{classroomqa}

The splitting is not produced by \(\langle F\rangle\) in isolation; it depends on couplings that transmit the breaking to the visible fields. If those couplings are weak, the visible particles remain recognizable and the splittings are small compared with grand-unification or Planck scales. If they were strong, the analogue would be an atom so distorted by an enormous magnetic field that it became a long, unrecognizable object.

\subsection{The Goldstino}

For an ordinary broken generator, \(J|0\rangle\) gives a bosonic Goldstone state. A supercharge is fermionic and changes bosonic states into fermionic ones. Since the vacuum is counted as bosonic, a broken supercharge creates a fermionic zero-energy state:
\begin{equation}
Q_\alpha|0\rangle
\sim
|\widetilde G_\alpha,\mathbf k=0\rangle.
\label{eq:l9-goldstino}
\end{equation}
Its propagating continuation is a massless Weyl fermion, the Goldstone fermion or goldstino.

\begin{classroomqa}
\audiencequestion{If a supersymmetry transformation does not leave the vacuum fixed, what new state does it make? \lecturetimestamp{00:53:24}}
\lecturerresponse{A massless fermion. The ordinary Goldstone theorem pattern remains, but a fermionic generator turns the bosonic vacuum into the goldstino.}
\end{classroomqa}

Rigid spontaneous supersymmetry breaking therefore has two unavoidable signatures:
\begin{equation}
\boxed{
V_{\min}>0
\qquad\text{and}\qquad
\text{a massless goldstino}.
}
\label{eq:l9-rigid-breaking-signatures}
\end{equation}

\section{The Hidden Sector}

It is useful to call the atom and the ferromagnet two sectors. The ferromagnet develops the order parameter; the atom feels only the field transmitted from it. When the coupling is weak, manipulating the atom scarcely excites a magnon and gives little direct access to the magnetic medium. The ferromagnet is then hidden, though its presence is inferred from small spectral splittings.

Supersymmetric model building uses the same language. Additional fields in a hidden sector develop nonzero \(F\)- or \(D\)-term expectation values. Interactions called the mediation mechanism communicate that breaking to the visible sector and generate soft masses for superpartners. Weak transmission explains how the visible theory can remain approximately supersymmetric at a much higher fundamental scale while showing no exact low-energy boson--fermion degeneracy.

The construction is not automatic. Extra fields, a stable breaking vacuum, and a controlled mediation mechanism are needed. This is why the lecture describes realistic supersymmetry breaking as complicated and somewhat contrived rather than as a single inevitable consequence of the minimal model.

\section{Gravity Changes Both Signatures}

\subsection{Vacuum Energy Becomes a Source}

Without gravity, shifting every energy by a constant does not change nongravitational experiments. With gravity, a vacuum energy density enters Einstein's equation as a cosmological constant. A positive \(V_{\min}\) drives accelerated expansion. If its scale were comparable to the particle-physics scale needed for superpartner masses, the resulting curvature would be enormously larger than observed.

\begin{classroomqa}
\audiencequestion{What does a positive nonzero vacuum energy do to the world once gravity is included? \lecturetimestamp{01:00:40}}
\lecturerresponse{It acts as a cosmological constant and drives accelerated expansion. At ordinary particle-physics scales the effect would be intolerably large.}
\end{classroomqa}

Local supersymmetry, or supergravity, alters the positive-square formula. For orientation, the standard \(N=1\) \(F\)-term potential is
\begin{equation}
\begin{aligned}
V_F={}&e^{K/M_P^2}
\left(
K^{i\bar j}D_iW\,D_{\bar j}\overline W
-\frac{3|W|^2}{M_P^2}
\right),\\
D_iW={}&\partial_iW+\frac{\partial_iK}{M_P^2}W.
\end{aligned}
\label{eq:l9-sugra-potential}
\end{equation}
The negative term means that a broken-supersymmetry contribution can in principle be balanced to leave a small net cosmological constant.

\editorialnote{The lecture intentionally gives only a schematic correction and explicitly questions the placement of the Planck mass. Equation~\eqref{eq:l9-sugra-potential} records the conventional formula for chiral multiplets with Kähler potential \(K\); gauge \(D\)-terms would add a further nonnegative contribution. It does not explain the observed tiny value, which remains the cosmological-constant problem.}

\subsection{The Super-Higgs Mechanism}

The second gravitational effect is easiest to understand by counting polarizations. A massive spin-\(j\) particle has \(2j+1\) spin states in its rest frame.

\begin{classroomqa}
\audiencequestion{How many spin states does a massive spin-\(j\) particle have? \lecturetimestamp{01:05:16}}
\lecturerresponse{\(2j+1\). In particular a massive spin-one particle has three states, whereas a massless gauge boson has only helicities \(+1\) and \(-1\).}
\end{classroomqa}

The distinction is kinematic. A massive particle can be boosted to rest and then sent in a new direction while its spin orientation is retained, so all projections must exist. A massless particle cannot be overtaken or brought to rest. It may therefore carry only the two extremal helicities.

For the electroweak gauge field,
\begin{equation}
\underbrace{(+1,-1)}_{\text{massless vector}}
+
\underbrace{(0)}_{\text{Goldstone}}
\longrightarrow
\underbrace{(+1,0,-1)}_{\text{massive vector}}.
\label{eq:l9-higgs-count}
\end{equation}

\begin{classroomqa}
\audiencequestion{Where does the longitudinal polarization of a massive \(W\) or \(Z\) come from? \lecturetimestamp{01:08:50}}
\lecturerresponse{From the Goldstone boson. It is absorbed by the gauge field and completes the three-state massive spin-one representation.}
\end{classroomqa}

Gravity has a massless spin-two graviton. Its supersymmetric partner is the massless spin-\(\tfrac32\) gravitino.

\begin{classroomqa}
\audiencequestion{What is the spin of the graviton, and therefore of its superpartner? \lecturetimestamp{01:10:24}}
\lecturerresponse{The graviton has spin \(2\); the fermionic partner differs by one half and has spin \(3/2\).}
\end{classroomqa}

A massless gravitino has helicities \(\pm\tfrac32\), while a massive gravitino must also have \(\pm\tfrac12\). The goldstino supplies precisely those missing states:
\begin{equation}
\begin{gathered}
\underbrace{\left(+\frac32,-\frac32\right)}_{\text{massless gravitino}}
\oplus
\underbrace{\left(+\frac12,-\frac12\right)}_{\text{goldstino}}\\
\Downarrow\\
\underbrace{\left(+\frac32,+\frac12,-\frac12,-\frac32\right)}_{\text{massive gravitino}}.
\end{gathered}
\label{eq:l9-super-higgs}
\end{equation}
This is the super-Higgs mechanism. Once supersymmetry is local, the goldstino is absorbed and no longer appears as an independent massless particle. The resulting gravitino interacts gravitationally and is correspondingly difficult to produce.

\section{Changing Subjects: Why \texorpdfstring{\(SU(5)\)}{SU(5)}?}

The supersymmetry formalism now gives way to grand unification. The Standard Model gauge group is
\begin{equation}
G_{\mathrm{SM}}
=SU(3)_c\times SU(2)_L\times U(1)_Y.
\label{eq:l9-sm-group}
\end{equation}
Its three factors separately organize color, weak isospin, and weak hypercharge. The unification question is whether this product is the surviving subgroup of one larger simple group, so that apparently separate particle multiplets and gauge forces are pieces of a single structure. The minimal candidate developed here is \(SU(5)\).

\subsection{Group Elements and Generators}

For \(U\in SU(N)\),
\begin{equation}
U^\dagger U=\mathbf1,
\qquad
\det U=1.
\label{eq:l9-SUN-def}
\end{equation}
The first condition is unitary; the second is the meaning of special. Write an infinitesimal element as
\begin{equation}
U=\mathbf1+i\epsilon G+\mathcal O(\epsilon^2).
\label{eq:l9-infinitesimal-U}
\end{equation}
Unitarity requires
\begin{equation}
G^\dagger=G,
\end{equation}
and
\begin{equation}
\det e^{i\epsilon G}
=e^{i\epsilon\operatorname{tr}G}=1
\end{equation}
requires
\begin{equation}
\operatorname{tr}G=0.
\end{equation}
Thus the generators are traceless Hermitian matrices.

An \(N\times N\) Hermitian matrix contains \(N^2\) real parameters: \(N\) real diagonal entries and \(N(N-1)/2\) complex off-diagonal entries. Tracelessness removes one. Therefore
\begin{equation}
\dim SU(N)=N^2-1,
\qquad
\dim SU(5)=24.
\label{eq:l9-SU5-dimension}
\end{equation}
For comparison, \(SU(2)\) has three Pauli generators and \(SU(3)\) has eight color generators.

\section{The Standard Model Blocks Inside \texorpdfstring{\(SU(5)\)}{SU(5)}}

Split the five-dimensional space into a weak two-dimensional block and a color three-dimensional block:
\begin{equation}
\mathbf5=\mathbf2\oplus\mathbf3.
\end{equation}
The embedded generators have the forms
\begin{align}
T^a_L
&=
\begin{pmatrix}
\tfrac12\tau^a&0\\
0&0_{3\times3}
\end{pmatrix},
&& a=1,2,3,
\label{eq:l9-SU2-block}\\
T^A_c
&=
\begin{pmatrix}
0_{2\times2}&0\\
0&\tfrac12\lambda^A
\end{pmatrix},
&& A=1,\ldots,8.
\label{eq:l9-SU3-block}
\end{align}
The Pauli matrices \(\tau^a\) generate weak \(SU(2)\); the Gell-Mann matrices \(\lambda^A\) generate color \(SU(3)\). Because the blocks do not overlap, the two subgroups commute.

One additional diagonal generator is proportional to the identity within each block and therefore commutes with both. In the fundamental representation one convenient doubled-hypercharge convention is
\begin{equation}
T_Y^{(\mathbf5)}=
\operatorname{diag}
\left(
1,1,-\frac23,-\frac23,-\frac23
\right).
\label{eq:l9-hypercharge-generator}
\end{equation}
Its trace vanishes:
\begin{equation}
1+1-3\left(\frac23\right)=0.
\end{equation}
The conjugate matter representation has the opposite eigenvalues,
\begin{equation}
Y_{\overline{\mathbf5}}
=
\operatorname{diag}
\left(
-1,-1,\frac23,\frac23,\frac23
\right),
\label{eq:l9-fivebar-hypercharge}
\end{equation}
which are the values written for the matter column in the lecture. The overall sign and normalization are conventional; these entries use the Standard Model convention
\begin{equation}
Q=T_3+\frac{Y}{2}.
\label{eq:l9-charge}
\end{equation}
Equations~\eqref{eq:l9-SU2-block}--\eqref{eq:l9-hypercharge-generator} display
\begin{equation}
SU(3)_c\times SU(2)_L\times U(1)_Y
\subset SU(5).
\end{equation}

The three weak generators, eight color generators, and one hypercharge generator account for \(12\) of the \(24\) \(SU(5)\) generators.\editorialnote{The other twelve occupy the off-diagonal \(2\times3\) and \(3\times2\) blocks. Their gauge bosons convert quark-type and lepton-type components into one another; this is the structural origin of the baryon- and lepton-changing interactions developed in the continuation.}

\begin{classroomqa}
\audiencequestion{Which gauge particles correspond to the familiar subgroup generators? \lecturetimestamp{01:31:05}}
\lecturerresponse{The eight \(SU(3)\) generators give the eight gluons. The three \(SU(2)\) generators give the weak gauge fields, which after electroweak mixing become \(W^\pm\) and part of the \(Z\)-photon system; the \(U(1)\) field supplies the other neutral combination.}
\end{classroomqa}

\section{One Family in Left-Handed Language}

Focus on the first generation; the same pattern is repeated for the muon/charm/strange and tau/top/bottom generations. It is most economical to count every independent Weyl field as left-handed. Right-handed particles are represented by their left-handed charge-conjugate fields:
\begin{equation}
Q_L=(u_L,d_L),\quad
L_L=(\nu_{eL},e_L),\quad
u^c_L,\quad d^c_L,\quad e^c_L.
\end{equation}
Their multiplicities are
\begin{equation}
\underbrace{6}_{Q_L}
+
\underbrace{2}_{L_L}
+
\underbrace{3}_{u^c_L}
+
\underbrace{3}_{d^c_L}
+
\underbrace{1}_{e^c_L}
=15.
\label{eq:l9-fifteen-fields}
\end{equation}

\editorialnote{The minimal Standard Model contains no independent right-handed-neutrino Weyl field, so the count is \(15\). This counting alone does not assert that the neutrino is a Majorana particle. Adding a singlet \(\nu^c_L\) makes \(16\) fields and points naturally toward \(SO(10)\), beyond the present \(SU(5)\) construction.}

The five-dimensional multiplet must place colorless weak-doublet entries in the first two slots and colored weak singlets in the last three.

\begin{classroomqa}
\audiencequestion{Does the proposed column contain six fields rather than five? \lecturetimestamp{01:36:14}}
\lecturerresponse{No. The up quark written during the first pass is erased. The column is the neutrino, the electron, and three color copies of the anti-down quark: five entries in all.}
\end{classroomqa}

That gives
\begin{equation}
\overline{\mathbf5}
=
\begin{pmatrix}
\nu_{eL}\\
e_L\\
d^c_{rL}\\
d^c_{gL}\\
d^c_{bL}
\end{pmatrix}.
\label{eq:l9-fivebar}
\end{equation}
\editorialnote{The lecture initially calls this simply a five-dimensional multiplet. In standard all-left-handed \(SU(5)\) notation, its color and hypercharge assignments identify it as the conjugate representation \(\overline{\mathbf5}\).}
The first two entries transform under \(SU(2)_L\) but not color. The last three transform as an anti-triplet under color but are weak singlets.

\begin{classroomqa}
\audiencequestion{Which quarks carry color but no weak \(SU(2)_L\) quantum number, once every field is written left-handed? \lecturetimestamp{01:40:13}}
\lecturerresponse{The charge conjugates of the right-handed quarks. For this five-component multiplet, the charge test selects the three anti-down fields \(d^c_{r,g,b}\).}
\end{classroomqa}

The electric charges in Eq.~\eqref{eq:l9-fivebar} sum to zero:
\begin{equation}
0-1+\frac13+\frac13+\frac13=0.
\label{eq:l9-fivebar-charge-sum}
\end{equation}
That is required because the electric-charge operator is a traceless combination of \(SU(5)\) generators.

\begin{classroomqa}
\audiencequestion{Why do anti-down quarks, rather than anti-up quarks, occupy the three color slots? \lecturetimestamp{01:43:38}}
\lecturerresponse{The neutrino and electron contribute total charge \(-1\). Three anti-down colors contribute \(+1\), making the trace zero. Three anti-up fields would have the wrong total charge.}
\end{classroomqa}

The remaining ten left-handed fields are
\begin{equation}
Q_L\;(6)
\oplus
u^c_L\;(3)
\oplus
e^c_L\;(1),
\end{equation}
and they fit exactly into the antisymmetric rank-two representation
\begin{equation}
\mathbf{10}\cong\wedge^2\mathbf5,
\qquad
\dim\wedge^2\mathbf5=\binom52=10.
\label{eq:l9-ten-preview}
\end{equation}
Their explicit placement is the next lecture's task. The numerical fit
\begin{equation}
\overline{\mathbf5}\oplus\mathbf{10}
\end{equation}
for every Standard Model family is the first striking evidence that the \(SU(5)\) idea is more than an arbitrary relabeling.

\begin{classroomqa}
\audiencequestion{If the five and ten are different representations, does that mean their particles cannot interact? \lecturetimestamp{01:46:14}}
\lecturerresponse{No. A group transformation acts within each irreducible representation, so it does not directly rotate a member of the five into a member of the ten. The two sets can nevertheless interact by emitting and absorbing common gauge bosons, and Higgs couplings can connect their fields as well.}
\end{classroomqa}

\section{Perspective}

The lecture has followed one pattern across very different scales. Rotational invariance makes atomic levels degenerate; a magnetic background splits them. A ferromagnet produces such a background while leaving the underlying laws symmetric, and its gapless magnons reveal that the asymmetry resides in the vacuum. In a chiral supersymmetric theory, the auxiliary expectation value plays the role of the order parameter:
\begin{equation}
\begin{aligned}
\langle F\rangle=0
&\quad\Longleftrightarrow\quad
\text{unbroken supersymmetry},\\
\langle F\rangle\neq0
&\quad\Longrightarrow\quad
\text{spontaneous breaking}.
\end{aligned}
\end{equation}
The broken vacuum carries positive energy in rigid supersymmetry, splits supermultiplets, and creates a goldstino. Weak mediation motivates a hidden sector; gravity turns vacuum energy into a cosmological problem and absorbs the goldstino into a massive gravitino.

The closing \(SU(5)\) construction begins a new question rather than completing it. Traceless Hermitian \(5\times5\) generators contain the Standard Model blocks, and the first five matter fields fit exactly as
\((\nu_e,e^-,d^c_r,d^c_g,d^c_b)\). The remaining ten are waiting in the antisymmetric representation. That unfinished fit is where grand unification properly begins.
